\section*{Bab \ref{ch:g12:complex} --- Bilangan Kompleks}

\begin{solution}{exo:g12:complex:1}
$(3-2\iu)(1+4\iu) = 3 + 12\iu - 2\iu - 8\iu^2 = 11 + 10\iu$.

$\dfrac{2+\iu}{1-\iu} = \dfrac{(2+\iu)(1+\iu)}{(1-\iu)(1+\iu)}
= \dfrac{1 + 3\iu}{2} = \dfrac12 + \dfrac32\iu$.

$\iu^{2027}$: karena $\iu^4 = 1$ dan $2027 = 4 \times 506 + 3$, maka
$\iu^{2027} = \iu^3 = -\iu$.
\end{solution}

\begin{solution}{exo:g12:complex:2}
\emph{(a)} $\Delta = 16 - 52 = -36$:
$z = \dfrac{4 \pm 6\iu}{2} = 2 \pm 3\iu$. Keduanya sepasang \omterm{def:g12:complex:conjugate}{konjugat}; hasil kalinya
$= 4 + 9 = 13 = \frac ca$. \emph{(b)} $\Delta = 1 - 4 = -3$:
$z = \dfrac{-1 \pm \iu\sqrt3}{2}$ (yaitu akar pangkat tiga dari satuan, $j$ dan
$\conj j$); hasil kalinya $= \frac{1 + 3}{4} = 1$.
\end{solution}

\begin{solution}{exo:g12:complex:3}
$\abs{z_1} = \sqrt2$, \omterm{def:g12:complex:argument}{argumennya} $\frac{3\pi}{4}$ (kuadran kedua), jadi
$z_1 = \sqrt2\,\eu^{3\iu\pi/4}$.
Lalu $\abs{z_2} = 2$, \omterm{def:g12:complex:argument}{argumennya} $-\frac\pi6$, jadi
$z_2 = 2\,\eu^{-\iu\pi/6}$. Karena itu
\[
\frac{z_1}{z_2} = \frac{\sqrt2}{2}\,\eu^{\iu\left(\frac{3\pi}{4} + \frac{\pi}{6}\right)}
= \frac{\sqrt2}{2}\,\eu^{11\iu\pi/12}.
\]
Secara aljabar,
\[
\frac{z_1}{z_2} = \frac{-1+\iu}{\sqrt3 - \iu}
= \frac{(-1+\iu)(\sqrt3+\iu)}{4}
= \frac{-\sqrt3 - 1 + \iu(\sqrt3 - 1)}{4}.
\]
Dengan menyamakan \omterm{def:g12:complex:def}{bagian realnya}:
$\frac{\sqrt2}{2}\cos\frac{11\pi}{12} = \frac{-\sqrt3-1}{4}$, sehingga
\[
\cos\frac{11\pi}{12} = -\frac{\sqrt6 + \sqrt2}{4}.
\]
\end{solution}

\begin{solution}{exo:g12:complex:4}
\emph{(a)} $\abs{z - 2} = \abs{z - (-\iu)}$ berarti $M$ berjarak sama dari
$A(2, 0)$ dan $B(0, -1)$: itulah \emph{sumbu ruas} $[AB]$.

\emph{(b)} Berjarak $3$ dari titik $(1,1)$: itulah \emph{lingkaran} berpusat
$1 + \iu$ dan berjari-jari $3$.

\emph{(c)} Menurut \cref{prop:g12:complex:geometry}(4), $\frac{z-1}{z+1}$ yang murni
imajiner berarti \omterm{def:g11:vect:vector}{vektor} dari $A(1,0)$ ke $M$ dan dari $B(-1,0)$ ke $M$
saling tegak lurus: jadi $M$ memandang ruas $[AB]$ dengan sudut siku-siku. Himpunannya
adalah \emph{lingkaran berdiameter $[AB]$} (yaitu lingkaran satuan) \emph{tanpa
titik $A$ dan $B$} itu sendiri.
\end{solution}

\begin{solution}{exo:g12:complex:5}
De Moivre: $\cos3\theta + \iu\sin3\theta = (\cos\theta + \iu\sin\theta)^3$.
Dengan \omterm{def:g10:algebra:expand}{menjabarkannya} lewat teorema binomial dengan $c = \cos\theta$, $s = \sin\theta$:
\[
(c + \iu s)^3 = c^3 + 3\iu c^2 s - 3 c s^2 - \iu s^3
= (c^3 - 3cs^2) + \iu(3c^2 s - s^3).
\]
Dengan menyamakan bagiannya dan memakai $s^2 = 1 - c^2$, $c^2 = 1 - s^2$:
\[
\cos3\theta = 4\cos^3\theta - 3\cos\theta,
\qquad
\sin3\theta = 3\sin\theta - 4\sin^3\theta .
\]
\end{solution}

\begin{solution}{exo:g12:complex:6}
\[
\cos^3\theta = \left(\frac{\eu^{\iu\theta} + \eu^{-\iu\theta}}{2}\right)^{\!3}
= \frac{\eu^{3\iu\theta} + 3\eu^{\iu\theta} + 3\eu^{-\iu\theta} + \eu^{-3\iu\theta}}{8}
= \frac{\cos3\theta + 3\cos\theta}{4}.
\]
Karena itu
\[
\int_0^{\pi/2}\cos^3\theta\,\dd\theta
= \frac14\left[\frac{\sin3\theta}{3} + 3\sin\theta\right]_0^{\pi/2}
= \frac14\left(-\frac13 + 3\right) = \frac{2}{3}.
\]
(Selaras dengan $W_3 = \frac23$ pada \cref{exo:g12:integ:9}, sebab
$\int_0^{\pi/2}\cos^3 = \int_0^{\pi/2}\sin^3$ menurut kesetangkupan
$\theta \mapsto \frac\pi2 - \theta$.)
\end{solution}

\begin{solution}{exo:g12:complex:7}
$c = a + \eu^{\pm\iu\pi/3}(b - a)$ dengan $b - a = 2 + \iu$ dan
$\eu^{\pm\iu\pi/3} = \frac12 \pm \iu\frac{\sqrt3}{2}$:
\[
\eu^{\iu\pi/3}(2+\iu) = \left(\tfrac12 + \iu\tfrac{\sqrt3}{2}\right)(2+\iu)
= \left(1 - \tfrac{\sqrt3}{2}\right) + \iu\left(\tfrac12 + \sqrt3\right),
\]
jadi $c_1 = \left(2 - \frac{\sqrt3}{2}\right) +
\iu\left(\frac32 + \sqrt3\right)$, dan serupa dengan sudut
$-\frac\pi3$:
$c_2 = \left(2 + \frac{\sqrt3}{2}\right) + \iu\left(\frac32 - \sqrt3\right)$.
\end{solution}

\begin{solution}{exo:g12:complex:8}
$-4 = 4\,\eu^{\iu\pi}$, sehingga penyelesaian $z^4 = -4$ adalah
\[
z_k = \sqrt2\,\eu^{\iu\left(\frac\pi4 + \frac{k\pi}{2}\right)},
\quad k = 0, 1, 2, 3,
\]
\emph{yaitu} $1 + \iu$, $-1 + \iu$, $-1 - \iu$, $1 - \iu$: itulah titik sudut
sebuah persegi berjari-jari luar $\sqrt 2$. Dengan mengelompokkan \omterm{def:g11:quad:discriminant}{akar} yang berkonjugat,
\[
z^4 + 4 = \bigl(z - (1{+}\iu)\bigr)\bigl(z - (1{-}\iu)\bigr)
\bigl(z + (1{-}\iu)\bigr)\bigl(z + (1{+}\iu)\bigr)
= (z^2 - 2z + 2)(z^2 + 2z + 2).
\]
\end{solution}

\begin{solution}{exo:g12:complex:9}
Dengan $q = \eu^{\iu\theta} \neq 1$:
\[
\sum_{k=0}^{n} \eu^{\iu k\theta} = \frac{\eu^{\iu(n+1)\theta} - 1}{\eu^{\iu\theta} - 1}
= \frac{\eu^{\iu(n+1)\theta/2}}{\eu^{\iu\theta/2}}
\cdot\frac{\eu^{\iu(n+1)\theta/2} - \eu^{-\iu(n+1)\theta/2}}{\eu^{\iu\theta/2} - \eu^{-\iu\theta/2}}
= \eu^{\iu n\theta/2}\,
\frac{\sin\frac{(n+1)\theta}{2}}{\sin\frac{\theta}{2}},
\]
dengan memakai $\eu^{\iu\alpha} - \eu^{-\iu\alpha} = 2\iu\sin\alpha$ (yaitu
\emph{pemfaktoran setengah sudut}). Dengan mengambil \omterm{def:g12:complex:def}{bagian realnya}:
\[
S = \frac{\sin\frac{(n+1)\theta}{2}}{\sin\frac{\theta}{2}}\,
\cos\frac{n\theta}{2}.
\]
\end{solution}

\begin{solution}{exo:g12:complex:10}
\emph{1.} $\omega$ adalah \omterm{def:g12:complex:rootsofunity}{akar satuan} ke-$5$ yang berbeda dari $1$, sehingga jumlah
kelima \omterm{def:g11:quad:discriminant}{akarnya} bernilai nol (\cref{thm:g12:complex:rootsofunity}), dan itu terbaca
$1 + \omega + \omega^2 + \omega^3 + \omega^4 = 0$.

\emph{2.} $u + v = \omega + \omega^2 + \omega^3 + \omega^4 = -1$ menurut butir 1.
Untuk hasil kalinya, dengan memakai $\omega^5 = 1$:
\[
uv = (\omega + \omega^4)(\omega^2 + \omega^3)
= \omega^3 + \omega^4 + \omega^6 + \omega^7
= \omega^3 + \omega^4 + \omega + \omega^2 = -1 .
\]
Jadi $u, v$ adalah \omterm{def:g11:quad:discriminant}{akar} $X^2 + X - 1 = 0$, yaitu
$\left\{\frac{-1+\sqrt5}{2}, \frac{-1-\sqrt5}{2}\right\}$. Sekarang
$u = \omega + \conj\omega = 2\cos\frac{2\pi}{5} > 0$ (sebab
$\frac{2\pi}{5} < \frac{\pi}{2}$), sehingga $u = \frac{-1 + \sqrt5}{2}$.

\emph{3.} $\cos\frac{2\pi}{5} = \frac u2 = \frac{\sqrt5 - 1}{4}$.
\end{solution}

\begin{solution}{pb:g12:complex:1}
\textbf{1.} $(2 + \iu)(3 - \iu) = 6 - 2\iu + 3\iu + 1 =
7 + \iu$. \quad $\frac{1 + \iu}{1 - \iu} =
\frac{(1 + \iu)^2}{2} = \frac{2\iu}{2} = \iu$. \quad
$\abs{3 + 4\iu} = 5$.

\textbf{2.} $1 + \iu = \sqrt2\,\eu^{\iu\pi/4}$;
$-2 = 2\eu^{\iu\pi}$;
$1 - \iu\sqrt3 = 2\eu^{-\iu\pi/3}$.

\textbf{3.} $(1 + \iu)^8 = \left(\sqrt2\right)^8
\eu^{8\iu\pi/4} = 16\,\eu^{2\iu\pi} = 16$.

\textbf{4.} $z^2 = \eu^{\iu\pi/2}$, jadi
$z = \pm\eu^{\iu\pi/4} = \pm\frac{\sqrt2}{2}(1 + \iu)$. Lalu
$\Delta = 4 - 20 = -16$, jadi $z = 1 \pm 2\iu$.

\textbf{5.} Pemetaan $z \mapsto \iu z$: putaran seperempat lingkaran
terhadap titik asalnya. Lalu $\abs{z - (1 + \iu)} = 2$: lingkaran
berpusat $1 + \iu$ dan berjari-jari $2$.

\textbf{6.} $\left(\omega^k\right)^5 = \eu^{2\iu\pi k} = 1$,
dan kelima titik $\eu^{2\iu k\pi/5}$ itu berbeda: jadi itulah kelima
penyelesaian $z^5 = 1$. Semuanya duduk pada lingkaran satuan dengan
sudut yang sama, $72^\circ$: yaitu segi lima beraturan dengan satu titik sudut di
$1$.

\textbf{7.} $S = 1 + \omega + \dots + \omega^4$ memenuhi
$(1 - \omega)S = 1 - \omega^5 = 0$, dan $\omega \neq 1$, jadi
$S = 0$. (Titik seimbang segi limanya adalah pusatnya.)

\textbf{8.} $\omega^4 = \overline\omega$ dan
$\omega^3 = \overline{\omega^2}$: dengan memasangkan \omterm{def:g12:complex:conjugate}{konjugatnya}, \omterm{def:g12:complex:def}{bagian
real} jumlahnya terbaca
$1 + 2\cos\frac{2\pi}{5} + 2\cos\frac{4\pi}{5} = 0$.

\textbf{9.} $\cos\frac{4\pi}{5} = 2c^2 - 1$, sehingga
$1 + 2c + 4c^2 - 2 = 0$, yaitu $4c^2 + 2c - 1 = 0$ dan
$c = \frac{-1 + \sqrt5}{4}$ (\omterm{def:g11:quad:discriminant}{akar} yang positif, sebab
$72^\circ$ sudut lancip):
$\cos 72^\circ = \frac{\sqrt5 - 1}{4}$.

\textbf{10.} Secara numeris $c \approx 0.30902$, cocok dengan
kalkulatornya. Lalu $\frac{1}{2\varphi} = \frac{1}{1 + \sqrt5} =
\frac{\sqrt5 - 1}{4}$ (dengan merasionalkannya): jadi $c = \frac{1}{2\varphi}$.
Segi lima itu keemasan sampai ke tulangnya: diagonalnya memotong
sisinya dengan nisbah $\varphi$ --- keserupaan diri yang tak berujung milik
bintang lima itu.

\textbf{11.} \omterm{def:g11:quad:discriminant}{Akar} kuadrat terlukis dengan jangka (lewat mesin
setengah lingkaran milik rata-rata ukur di jilid sekolah menengah
pertama), sehingga bilangan yang dibangun dari \omterm{def:g10:numbers:sets}{bilangan rasional} hanya
lewat akar kuadrat --- seperti $\cos 72^\circ$ --- terlukiskan; sedangkan akar
pangkat tiga tidak, dan $\cos 20^\circ$ memerlukannya: segi limanya tunduk,
pembagian tiganya melawan. Kaidah Gauss memutuskan setiap segi banyak
beraturan --- segi tujuh belasnya pada usia sembilan belas membuatnya memilih
matematika.

\textbf{12.} $-4 = 4\eu^{\iu\pi}$: \omterm{def:g11:quad:discriminant}{akar} keempatnya bermodulus
$\sqrt2$ dan berargumen
$\frac{\pi}{4} + k\frac{\pi}{2}$, yaitu keempat bilangan
$\pm 1 \pm \iu$. Dengan memasangkan \omterm{def:g12:complex:conjugate}{konjugatnya}:
$(z - 1 - \iu)(z - 1 + \iu) = z^2 - 2z + 2$ dan
$(z + 1 - \iu)(z + 1 + \iu) = z^2 + 2z + 2$, yaitu
pemfaktoran pada \cref{exo:g12:complex:8}.

\textbf{13.} $(\cos\theta + \iu\sin\theta)^3 = \cos^3\theta +
3\iu\cos^2\theta\sin\theta - 3\cos\theta\sin^2\theta -
\iu\sin^3\theta$; dan menurut de Moivre ini sama dengan
$\cos 3\theta + \iu \sin 3\theta$. \omterm{def:g12:complex:def}{Bagian realnya}:
$\cos 3\theta = \cos^3\theta - 3\cos\theta(1 - \cos^2\theta)
= 4\cos^3\theta - 3\cos\theta$ --- sedangkan \omterm{def:g12:complex:def}{bagian imajinernya}
menyerahkan $\sin 3\theta$ secara cuma-cuma.

\textbf{14.} Itulah \omterm{def:g12:complex:def}{bagian real}
$1 + \omega + \dots + \omega^4 = 0$: jumlahnya sama dengan $0$ ---
tepat pertanyaan 8, yang dilihat lewat
inti pada \cref{exo:g12:complex:9}.

\textbf{15.} $8\iu = 8\eu^{\iu\pi/2}$, jadi \omterm{def:g11:quad:discriminant}{akarnya}
$2\eu^{\iu\pi/6} = \sqrt3 + \iu$;
$2\eu^{5\iu\pi/6} = -\sqrt3 + \iu$;
$2\eu^{3\iu\pi/2} = -2\iu$.

\textbf{16.} $1 + \iu = \sqrt2\,\eu^{\iu\pi/4}$: jadi pemetaannya
memutar bidangnya sebesar $45^\circ$ lalu menskalakannya dengan $\sqrt2$.
Persegi satuannya menjadi persegi miring bersisi $\sqrt2$ ---
luasnya berlipat dua, terpuntir seperdelapan putaran.

\textbf{17.} $\abs{z_1 z_2}^2 = z_1 z_2 \overline{z_1 z_2} =
z_1\overline{z_1}\, z_2\overline{z_2} =
\abs{z_1}^2\abs{z_2}^2$: lalu ambil akar kuadratnya. Pangkat
$1 + \iu$: $1 + \iu$, $2\iu$, $-2 + 2\iu$, $-4$: masing-masing
$45^\circ$ lebih jauh berputar dan $\sqrt2$ kali lebih panjang --- sebuah
ulir yang membuka keluar lewat \omterm{def:g12:complex:modulus}{modulus}
$\sqrt2, 2, 2\sqrt2, 4, \dots$

\textbf{18.} Bilangan $\mathrm{j}$ adalah \omterm{def:g12:complex:rootsofunity}{akar satuan} ketiga
yang berbeda dari $1$, sehingga (lewat alasan pertanyaan 7 dengan $n = 3$)
$1 + \mathrm{j} + \mathrm{j}^2 = 0$. Kaidahnya pada segitiga
bakunya: $a + \mathrm{j}b + \mathrm{j}^2 c = 1 +
\mathrm{j}\cdot\mathrm{j} + \mathrm{j}^2 \cdot \mathrm{j}^2 =
1 + \mathrm{j}^2 + \mathrm{j}^4 = 1 + \mathrm{j}^2 +
\mathrm{j} = 0$: jadi terpenuhi, seperti yang dituntut segitiga sama sisi
$\left(1, \mathrm{j}, \mathrm{j}^2\right)$.

\textbf{19.} Atas $\C$, suku banyak real terbelah menjadi faktor
linear; \omterm{def:g11:quad:discriminant}{akarnya} yang tak real datang berpasangan \omterm{def:g12:complex:conjugate}{konjugat} (konjugatkanlah
\omterm{def:g10:algebra:equation}{persamaannya}), dan tiap pasangannya terkalikan menjadi bentuk kuadrat real
$z^2 - 2\operatorname{Re}(z_0)\,z + \abs{z_0}^2$: dari situlah
pemfaktoran realnya menjadi bagian berderajat $1$ dan berderajat $2$ ---
dan pertanyaan 12 melakukannya pada $z^4 + 4$.

\textbf{20.} Titik; putaran-dan-penskalaan; segi banyak beraturan;
identitas trigonometri lewat penjabaran; dan sebuah dunia tertutup
tempat setiap \omterm{def:g10:algebra:equation}{persamaan} berderajat $n$ memiliki $n$ \omterm{def:g11:quad:discriminant}{akarnya}. Koda:
pada pertanyaan 10, segi lima milik $\pi$ berjabat tangan dengan nisbah
emas --- dua tamu Euklides yang paling masyhur, yang akhirnya diperkenalkan
oleh jumlah lima anak panah yang sama dengan nol.

\end{solution}
